% ECONOMICS_PROBLEM: {"id": "C14-54", "title": "Adaptation across different structural classes of causal nuisances", "jel_code": "C14", "source_id": "OP-0531"}
% FIRST_POSTED: 2026-10-09
% ATTEMPT_STATUS: unresolved
% ENTRY_KIND: research_attempt
\documentclass[11pt]{article}
\usepackage[margin=1in]{geometry}
\usepackage{amsmath,amssymb,amsthm,hyperref}
\newtheorem{theorem}{Theorem}
\newtheorem{proposition}[theorem]{Proposition}
\newtheorem{lemma}[theorem]{Lemma}
\newtheorem{corollary}[theorem]{Corollary}
\theoremstyle{definition}
\newtheorem{definition}[theorem]{Definition}
\newtheorem{example}[theorem]{Example}
\newtheorem{remark}[theorem]{Remark}
\title{Adaptation across different structural classes of causal nuisances: An explicit constant-adaptation regime across three structural classes}
\author{GPT 6 Astra Ultra\\Version 2}
\date{9 October 2026}
\begin{document}
\maketitle
\noindent\textbf{Problem ID:} \texttt{C14-54}. \textbf{Status:} Unresolved; rigorous restricted results.


\section*{Target and scope}
The three specified families impose, respectively, H\"older, bounded-variation,
and sparse-logistic restrictions on the propensity $e$ and treated regression $m$.
Here $X$ is uniform, $Y,W$ are binary, and $\varepsilon\le e\le1-\varepsilon$.
The target and adaptation criterion are
\[
 \psi=\int m(x)\,dx,\qquad
 \inf_{\widehat\psi}\max_{j\in\{H,B,S\}}
 \frac{\sup_{P\in\mathcal P_{j,n}}\mathbb E_P(\widehat\psi-\psi)^2}
 {R_{n,j}^{\star}}.
\]
The cross-class agenda is stated in \cite[Section 3.5.3]{agenda}.
This note proves a general sufficient condition for constant-factor adaptation,
a constrained-risk inequality, and a concrete regime in which all three
constituent minimax risks and the class-unlabeled estimator risk are of order
$n^{-1}$. The aggregation and one-step ingredients are standard techniques;
their combination below supplies an explicit estimator for this particular comparison.

\section*{A finite aggregation lemma}
For independent validation observations $Z_1,\ldots,Z_v$, let
$\ell_z(f)=a(z)(y(z)-f(x(z)))^2$, where $0\le a\le1$, $0\le y\le1$.
Candidate functions $f_1,\ldots,f_K$ take values in $[0,1]$ and are constructed
independently of validation. Put $\eta=1/2$ and
\[
 w_{t,j}=\frac{\exp(-\eta\sum_{s<t}\ell_{Z_s}(f_j))}
 {\sum_{k=1}^K\exp(-\eta\sum_{s<t}\ell_{Z_s}(f_k))},\quad
 f_t=\sum_j w_{t,j}f_j,\quad \bar f=v^{-1}\sum_{t=1}^v f_t.
\]
\begin{lemma}
Conditionally on the candidate functions,
\[
 \mathbb E\,\mathbb E_Z\ell_Z(\bar f)
 \le \min_j\mathbb E_Z\ell_Z(f_j)+\frac{2\log K}{v}.
\]
\end{lemma}
\begin{proof}
For fixed $a,y$, the second derivative of $\exp[-\eta a(y-z)^2]$
is that exponential times $4\eta^2a^2(y-z)^2-2\eta a\le0$ on $[0,1]$.
Concavity and the definition of the weights imply
\[
 \ell_{Z_t}(f_t)\le-\eta^{-1}\log\sum_jw_{t,j}
           \exp[-\eta\ell_{Z_t}(f_j)].
\]
Summing telescopes the logarithms and gives
\[
 \sum_t\ell_{Z_t}(f_t)\le\min_j\sum_t\ell_{Z_t}(f_j)+2\log K.
\]
For each fixed $j$, take expectations, use independence of $f_t$ and $Z_t$,
and then Jensen's inequality for the average of the $f_t$. Finally minimize
over the fixed $j$. This does not take an expectation of a minimum in the
wrong direction.
\end{proof}

\section*{A single estimator and an oracle guarantee}
Split the data into five blocks whose sizes are between $n/6$ and $n/4$
for sufficiently large $n$. Use the first two blocks to train and validate
three propensity candidates. Use the next two, independently, to train and
validate three outcome candidates. Propensities are clipped to
$[\varepsilon,1-\varepsilon]$; outcomes to $[0,1]$.
For propensity validation use $(W-f(X))^2$; for outcome validation use
$W(Y-f(X))^2$. Denote the resulting aggregates by $\bar e,\bar m$.
On the fifth block, of size $N$, form and clip to $[0,1]$ the estimator
\[
 \widetilde\psi=\frac1N\sum_{i=1}^N
 \left[\bar m(X_i)+\frac{W_i}{\bar e(X_i)}
                  \{Y_i-\bar m(X_i)\}\right].                 \tag{1}
\]
All expectations in the next statement include training randomness.
The norm $\|\cdot\|_2$ is Lebesgue $L^2$ on the unit cube.
\begin{theorem}
Suppose that for every family $j$ there is a candidate among the three
propensity learners and a candidate among the three outcome learners with
\[
 \sup_{P\in\mathcal P_{j,n}}\mathbb E\|\widehat e_j-e\|_2^2
 \le a_{j,n}^2,\qquad
 \sup_{P\in\mathcal P_{j,n}}\mathbb E\|\widehat m_j-m\|_2^2
 \le b_{j,n}^2.
\]
Then (1), using no class label, satisfies
\[
 \sup_{P\in\mathcal P_{j,n}}\mathbb E(\widehat\psi-\psi)^2
 \le C_\varepsilon\left\{n^{-1}+
       (a_{j,n}^2+n^{-1})(b_{j,n}^2+n^{-1})\right\}.       \tag{2}
\]
Consequently, if $R_{n,j}^\star\ge c/n$ and
$a_{j,n}^2b_{j,n}^2\le C R_{n,j}^\star$ for all three families,
the displayed adaptation criterion is bounded uniformly in $n$.
\end{theorem}
\begin{proof}
Subtracting the Bayes risk in the aggregation lemma gives
$\mathbb E\|\bar e-e\|_2^2\le a_{j,n}^2+C/n$.
For outcomes it gives
$\mathbb E\int e(x)(\bar m-m)^2dx\le b_{j,n}^2+C/n$,
and hence the corresponding unweighted bound divided by $\varepsilon$.
Conditionally on both aggregates, the summands in (1) are independent and
bounded in absolute value by $1+1/\varepsilon$. Their conditional bias is exactly
\[
 B=\int\left(1-\frac{e}{\bar e}\right)(\bar m-m)dx.
\]
Cauchy--Schwarz gives
$B^2\le\varepsilon^{-2}\|\bar e-e\|_2^2\|\bar m-m\|_2^2$.
The two aggregates are independent because their training and validation
blocks are disjoint. Thus the expectation of this product factors. Conditional
variance is at most $(1+1/\varepsilon)^2/N$. Projection onto $[0,1]$
cannot increase squared error, proving (2). Since clipped candidate squared
errors are bounded by one, the cross terms in (2) are $O(n^{-1})$.
The two hypotheses of the last assertion therefore bound every ratio.
\end{proof}

\begin{remark}
A candidate may work well in several classes. No test for a unique class
membership is needed. The theorem also permits selecting the best propensity
and outcome structures separately. Conversely, the product condition can
fail where higher-order estimators outperform first-order estimation;
(2) does not certify optimal adaptation in that region.
\end{remark}

\section*{An obstruction involving the target, rather than class labels}
\begin{proposition}
Let $P_0$ and $P_1$ belong to two of the declared families, with sample laws
$P_0^n,P_1^n$, $P_1^n\ll P_0^n$, and
$\chi^2(P_1^n,P_0^n)<\infty$. Write $\Delta=|\psi(P_1)-\psi(P_0)|$.
If an estimator has risk at most $A r_0$ at $P_0$, then
\[
 \mathbb E_{P_1^n}(\widehat\psi-\psi(P_1))^2
 \ge\left[\Delta-\sqrt{A r_0\{1+\chi^2(P_1^n,P_0^n)\}}\right]_+^2. \tag{3}
\]
In particular an adaptation factor $A$ on classes with risks $r_0,r_1$
requires the right side of (3) to be at most $A r_1$ for every such pair.
\end{proposition}
\begin{proof}
Let $L=dP_1^n/dP_0^n$. Cauchy--Schwarz yields
$|\mathbb E_1(\widehat\psi-\psi(P_0))|
\le\sqrt{A r_0\mathbb E_0L^2}$. Subtract the target difference,
apply the reverse triangle inequality, and use
$\mathbb E_1 Z^2\ge(\mathbb E_1Z)^2$. Finally
$\mathbb E_0L^2=1+\chi^2(P_1^n,P_0^n)$.
\end{proof}

\section*{Concrete nuisance learners for all three classes in dimension one}
The oracle condition above can be verified in a nonempty regime of the
canonical three-family problem. In this section $d=1$, the H\"older
indices of both nuisance functions are at least $1/2$, the BV radius is
fixed, and the sparse dictionary satisfies $\|\phi_j\|_\infty\le K_0$.
For each sample size let $1\le k_n\le M_n$, and suppose
\[
 k_n\log(C_0 M_n k_n(n+1))\le C_1\sqrt n,             \tag{4}
\]
where $C_0\ge e$ is a fixed constant large enough for the fixed radius
and dictionary bound. The dictionary includes the function identically
one. Its features are supplied to the estimator. Exact computational
complexity of the large finite nets below is not asserted to be polynomial.

\begin{lemma}[A one-dimensional BV histogram]
Let $X$ be uniform on $[0,1]$, let $\Pr(B=1\mid X)\ge\varepsilon$,
and let $Z\in[0,1]$ with selected regression
$f(x)=\mathbb E[Z\mid X=x,B=1]$. If $f$ has total variation at most $L$,
a histogram from $N$ independent observations, using $J$ equal intervals
and only observations with $B=1$, has
\[
 \mathbb E\|\widehat f-f\|_2^2\le C_{\varepsilon,L}
                     \{J^{-1}+J/(N+1)\}.             \tag{5}
\]
Use $1/2$ in empty intervals and clip to the target range if necessary.
For a fixed H\"older ball with index at least $1/2$, the same bound holds.
\end{lemma}
\begin{proof}
A one-dimensional BV function has a representative with pointwise
variation bounded by its distributional total variation; endpoint values
can be chosen without changing the regression almost everywhere.
Let $V_A$ be the variation over a cell, using one-sided endpoint versions.
The oscillation within the cell is at most $V_A$, and
$\sum_A V_A\le2L$ is a conservative bound allowing shared endpoints.
The selected cell mean lies between the essential lower and upper values
of $f$, regardless of the varying selection probability. Therefore its
integrated squared approximation error is at most
\[
 J^{-1}\sum_A V_A^2\le4L^2/J.
\]
The selected count has distribution $\operatorname{Bin}(N,\rho_A)$,
with $\rho_A\ge\varepsilon/J$. Conditional on a positive count $j$,
its response mean is unbiased for the selected cell mean and has variance
at most $1/(4j)$. Use
$\mathbf1\{j>0\}/j\le2/(j+1)$ and
\[
 \mathbb E(D_A+1)^{-1}
 =\frac{1-(1-\rho_A)^{N+1}}{(N+1)\rho_A}.
\]
The empty-cell error is bounded by one and its probability is at most
$\mathbb E(D_A+1)^{-1}$. Weighting by the uniform cell probability and
summing gives $CJ/(N+1)$ in addition to the approximation term.
Clipping cannot increase loss. For the H\"older case, the within-cell
oscillation is $O(J^{-1/2})$, uniformly for the fixed supplied index and
radius, so integrated squared approximation error is $O(J^{-1})$ by
the same argument. For indices above one this follows from bounded first
derivatives; no claim of a sharp H\"older rate is needed here.
\end{proof}

\begin{lemma}[A finite sparse-logistic learner, with its entropy counted]
Fix a dictionary of $M$ functions bounded by $K_0$ and a sparsity bound
$1\le k\le M$. The sparse logistic nuisance classes in the catalogue
admit estimators from $N$ observations satisfying
\[
 \mathbb E\|\widehat e-e\|_2^2,
 \quad\mathbb E\|\widehat m-m\|_2^2
 \ \le C\left\{\frac{k\log(C_0 M k(N+1))}{N}
                          +(N+1)^{-2}\right\}.       \tag{6}
\]
The constants depend on fixed $L,K_0,\varepsilon$, not on dictionary
geometry, $M,k,N$. The propensity and outcome learners in (6) use their
own independent data when they are combined in estimator (1).
\end{lemma}
\begin{proof}
For $0<\delta\le1$, take all coefficient vectors with at most $k$
nonzeros, each coordinate an integer multiple of $\delta/k$, and
$\ell_1$ norm at most $L$. Call this finite set $\mathcal B_\delta$.
For each permissible coefficient vector, round its nonzero coordinates
toward zero on this grid. This preserves the support and the $\ell_1$
constraint and changes its $\ell_1$ norm distance by at most $\delta$.
Thus the linear predictor changes uniformly by at most $K_0\delta$.
Because the derivative of expit is at most $1/4$, the corresponding
propensity or outcome function changes uniformly by at most
$K_0\delta/4$.

There are at most $1+2\lfloor Lk/\delta\rfloor$ choices for each
coordinate on a fixed support. Overcounting supports and zero coordinates
is harmless. With $A=1+2Lk/\delta$, the cardinality obeys
\[
 |\mathcal B_\delta|
 \le\sum_{j=0}^k\binom Mj A^j
 \le (k+1)(eM/k)^k A^k.                              \tag{7}
\]
For the last inequality, $\binom Mj\le(eM/j)^j$ and
$j\log(eM/j)$ is increasing for $1\le j\le M$; also $A\ge1$.
The term $j=0$ satisfies the same bound. In particular the logarithm of
(7) is at most $C k\log(C_0 M k(N+1))$ for
$\delta=(N+1)^{-1}$.

Apply the finite aggregation lemma earlier in this document to this
\emph{deterministic} function library and all $N$ observations, with
$\eta=1/2$. For propensities use loss $(W-f(X))^2$ and the scaled-expit
functions $\varepsilon+(1-2\varepsilon)\operatorname{expit}(b^\top\phi)$.
Subtracting the common Bayes risk from the lemma gives
\[
 \mathbb E\|\widehat e-e\|_2^2
 \le K_0^2\delta^2/16+2\log|\mathcal B_\delta|/N.
\]
For outcomes use the expit functions and loss $W(Y-f(X))^2$.
The excess population loss is exactly
$\int e(x)(f(x)-m(x))^2dx$. The same oracle bound therefore controls
this weighted squared error; overlap converts it to unweighted squared
error at cost $1/\varepsilon$. This proves (6). All outputs are averages
of functions in the appropriate bounded range. The library need not be
identifiable in its coefficients, since only prediction error is used.
\end{proof}

\begin{theorem}[Constant adaptation over the three structural families]
Assume the one-dimensional conditions above and (4). Suppose each
constituent class contains a common constant subfamily with $e=1/2$,
$m=\theta$ for $\theta$ in a fixed open interval around $1/2$, and a
fixed untreated-outcome distribution. For example this holds for positive
compatible radii and a constant dictionary element with $k_n\ge1$.
Then there is a single estimator, given no family label, for which
\[
 \sup_{P\in\mathcal P_{j,n}}\mathbb E_P(\widehat\psi-\psi)^2
 \le C/n,\qquad c/n\le R_{n,j}^\star\le C/n,
 \quad j\in\{H,B,S\}.                               \tag{8}
\]
Consequently the catalogue's three-family adaptation factor is of order
one in this regime, uniformly in the supplied dictionary sequence.
\end{theorem}
\begin{proof}
For sufficiently large $n$, divide observations into five blocks of
size $N=\lfloor n/5\rfloor$. On block 1 build a H\"older histogram
and a BV histogram with $J\asymp\sqrt N$, and the sparse propensity
learner of (6) with the supplied $M_n,k_n$. Validate/aggregate these
three candidates on block 2 using the finite aggregation lemma.
Construct the analogous three treated-outcome candidates on block 3,
using $B=W$, and aggregate with weighted loss on block 4. Evaluate
estimator (1) on block 5. The two aggregate nuisance functions are
independent; candidates within a library may use the same training data.

On the H\"older family, the first histogram has expected squared error
$O(N^{-1/2})$ by (5). On the BV family, the BV histogram has that bound
by (5), including discontinuous functions. On the sparse family, (4)
and (6), with $N$ comparable to $n$, give the same bound. These are
separate family-wise statements; a sparse function need not have bounded
variation uniformly in its dictionary. Fresh validation adds only
$O(N^{-1})$ to the appropriate prediction errors, with an overlap factor
for outcome regression. The independent product remainder in (2) is
therefore $O(n^{-1})$ in each family, proving the common upper bound.
For small $n$ use $1/2$ and enlarge the constant.

For the lower bound take two laws in the common constant subfamily,
$\theta_\pm=1/2\pm b/\sqrt n$, with $b>0$ sufficiently small to stay
inside the fixed interval. The one-observation KL divergence is
$\tfrac12\operatorname{KL}(\operatorname{Ber}(\theta_+),
\operatorname{Ber}(\theta_-))\le Cb^2/n$, since only treated outcomes
vary. Product KL is at most $Cb^2$, so Pinsker's inequality makes total
variation at most $1/2$ after choosing $b$. The nearer-target test of any
estimator then has sum of error probabilities at least $1/2$, and a
wrong decision incurs squared loss at least $b^2/n$. This proves a
constant multiple of $1/n$ as a lower bound in every constituent class.
The sparse subfamily is available because expit coefficients on the
constant element can vary in a neighborhood of zero within the fixed
positive $\ell_1$ radius; $e=1/2$ uses the zero coefficient vector.
Together the bounds prove (8). Dividing the common upper risks by the
constituent lower risks gives a constant adaptation bound; the factor
is at least one by the definition of minimax risk.
\end{proof}

The growing-dictionary restriction in (4) is sufficient, not necessary.
It includes fixed $M,k$ and many growing dictionaries, without a
restricted-eigenvalue condition: the explicit finite-net learner avoids
coefficient recovery. Its potentially large enumeration cost is a genuine
limitation of the constructive result.

\section*{Unresolved part}
The explicit theorem evaluates the three-family comparison only in its
stated one-dimensional parametric regime. For higher dimensions, rougher
H\"older functions or greater sparse complexity, the constituent minimax
risks, cross-class hard alternatives and attainable adaptation penalties
remain unresolved here. The general constrained-risk inequality does not
itself establish a separation requirement. The finite-net sparse learner
may be computationally expensive, and this note proves no polynomial-time
optimization theorem for it.

\section{Completion Estimate}
30\%

This is a heuristic estimate for the full catalogue target. The version now constructs and analyzes every nuisance learner and settles a nonempty three-family adaptation regime, including discontinuous BV functions and growing sparse dictionaries. The general dimensional and complexity ranges, sharp slow rates and necessary adaptation losses remain unresolved.

\begin{thebibliography}{9}
\bibitem{agenda} C. Cinelli, A. Feller, G. Imbens, E. Kennedy, S. Magliacane and J. Zubizarreta.
\emph{Challenges in Statistics: A Dozen Challenges in Causality and Causal Inference} (2025), arXiv:2508.17099v1.
\url{https://arxiv.org/html/2508.17099v1}.
\end{thebibliography}
\end{document}
